% A complete algebraic structure theorem for ordinary arms/products.
% This is not a Lefschetz realization of a general one-descent core.
\section{Lefschetz realizations and Hilbert growth of ordinary positive differences}
An $O$-sequence means the Hilbert function of a finite standard
graded quotient of a polynomial ring. Its standard monomials
give a finite monomial order ideal with the same degree counts.
Write $\mathcal D h=(\Delta h)^+$, with $h_{-1}=0$.
All algebras in the proof may be taken over $\mathbb Q$.

\begin{lemma}[Lifting a differentiable Hilbert function]
If $h_0=1$, $h$ is finite and unimodal, and $\mathcal D h$
is an $O$-sequence, then there exists a monomial Artinian
algebra with Hilbert function $h$ and a strong Lefschetz element.
\end{lemma}
\begin{proof}
Let $m$ be the first mode of $h$, and choose a finite monomial
order ideal $\mathcal M$ with degree counts $\mathcal D h$.
It has $h_m$ monomials, all of degree at most $m$.
List them in nondecreasing degree. Form a list of endpoints
with $h_j-h_{j+1}$ occurrences of $j$ for each $j\ge m$,
and list these in nonincreasing order. Pair the two lists,
assigning an endpoint $e(u)$ to each $u\in\mathcal M$.
Define the monomial basis
\[
 \mathcal B=\{u z^k:u\in\mathcal M,\ 0\le k\le e(u)-\deg u\}.
\]
This is an order ideal. Indeed, if $v\mid u$, then
$\deg v\le\deg u$ and $e(v)\ge e(u)$; every divisor
$v z^i$ of $u z^k$ therefore has total degree at most $e(v)$.
Division by $z$ also stays inside the displayed set.

The $z$ strings occupy intervals $[\deg u,e(u)]$. These
intervals are nested by construction. Their degree counts are
$h$: through $m$ the starts have counts $\mathcal D h$,
and after $m$ the prescribed endpoints remove precisely
$h_j-h_{j+1}$ strings at a time. At any two degrees the sets
of active nested strings are comparable by inclusion. Thus
multiplication by every power of $z$ has rank equal to the
smaller of the two graded dimensions. This proves the SLP.
\end{proof}

The general Hilbert-function equivalence behind this lemma is
classical; compare Harima--Wachi,
\emph{Generic initial ideals, graded Betti numbers and
$k$-Lefschetz properties}, Proposition~3 and the $k=1$ case
of Proposition~33 in the arXiv version
\url{https://arxiv.org/abs/0707.2247}.
The explicit nested-string construction above is included to
make the use of this background transparent. No priority is
claimed for this general realization criterion.

\begin{lemma}[A uniform window preserves existence of an SLP realization]
If $h$ is the Hilbert function of an SLP algebra, then
$[t]h$ has such a realization for every $t\ge1$.
In particular $\mathcal D([t]h)$ is an $O$-sequence.
\end{lemma}
\begin{proof}
Let $A$ realize $h$, with strong Lefschetz element $\ell$, and
put $B=A[y]/(y^t)$. Its Hilbert function is $b=[t]h$.
The quotient by $y+\ell$ is $A/(\ell^t)$ and has Hilbert
function
\[
                    \max(h_j-h_{j-t},0)
                         =\max(b_j-b_{j-1},0).
\]
Hence multiplication by $y+\ell$ has maximal rank between
consecutive degrees of $B$: it has the WLP. This implies
unimodality, since once a graded component of $B/(y+\ell)$
vanishes all later components vanish. Its Hilbert function
is $\mathcal D b$, so the lifting lemma gives an SLP
realization of $b$. We do not assert that the particular
tensor product $B$ always has the SLP.
\end{proof}

\begin{theorem}[Ordinary arms and two-arm products]
Let $q_X,q_Y$ be reversed formal arm polynomials with initial
runs of at least one one and no later parts equal to one.
The non-one parts may occur in arbitrary internal order.
Each of $q_X$ and $q_Xq_Y$, and its convolution with any
finite product of interval polynomials, has an SLP algebra
realization. In particular its positive differences form an
$O$-sequence.
\end{theorem}
\begin{proof}
For one arm the initial interval is a one-variable complete
intersection. The formal recurrence is
\[
                   q_{Zt}=[t]q_Z+z^{a},\qquad a=\operatorname{wt}Z+t.
\]
The window has degree $a-1$ and final coefficient one.
Appending one at rank $a$ preserves unimodality and does not
change its positive differences. The two lemmas therefore
give a full induction for single arms.

For two arms use induction on their combined skeleton length.
A rectangular factor is the window lemma. Otherwise orient
$a\ge b$. The established ordinary-product midpoint argument,
including its extra rank for equal nonrectangular factors,
says that the first mode $s$ of $q_Xq_Y$ satisfies $s<a$.
Deleting the last part $t$ of $X$ gives
\[
                        q_Xq_Y=[t]q_Zq_Y+z^a q_Y.
\]
Up to rank $s$, the positive differences are those of the
window term, whose $O$-sequence property follows by induction
and the window lemma. After $s$ they are zero by the already
proved ordinary-product unimodality. Thus $\mathcal D(q_Xq_Y)$
is the degree truncation of an $O$-sequence, again an
$O$-sequence. Apply the lifting lemma. Finally use the
window lemma repeatedly for additional interval factors.
\end{proof}

\begin{theorem}[Strict relative kernels]
Let $Y$ be increasing, every part smaller than $M\ge2$,
and $b=\operatorname{wt}Y$. For
$D_M(Y)=1+(S_M-1)F_Y$, the reversed kernel
\[
                 \delta_M(Y)=z^{b+M}D_M(Y)(1/z)
\]
has the equivalent expression
$\delta_M(Y)=[M]q_Y-z^Mq_Y+z^{b+M}$.
It has an SLP algebra realization, and
$\mathcal D\delta_M(Y)$ is an $O$-sequence.
\end{theorem}
\begin{proof}
Use the established unimodality and the append identity
$\delta_M(Zt)=[t]\delta_M(Z)+z^b[M-t]$, where $b=\operatorname{wt}(Zt)$.
If $b<M$, the kernel agrees with the cumulative coefficients
of $q_Y$ below $M$ and is nonincreasing from rank $b$ onward.
Its positive differences are exactly $q_Y$, already an
$O$-sequence by the ordinary-arm theorem.

If $b\ge M$, the previously proved joining inequality is
$\delta_M(Y)_b\le\delta_M(Y)_{b-1}$, and it is nonincreasing
afterward. The interval correction starts at $b$, so the
positive differences are the degree truncation of the
positive differences of $[t]\delta_M(Z)$. Induction and the
window lemma give the $O$-sequence property. The lifting
lemma completes both branches. The empty arm is the
initial interval $[M]$.
\end{proof}

\begin{corollary}[An adjacent negative-binomial envelope for positive differences]
For any polynomial $q$ in these theorems put $r=q_1$ and
$d=\mathcal Dq$. Then $d_0=1$, $d_1=r-1$, and
\[
                   (j+1)d_{j+1}\le(r-1+j)d_j\qquad(j\ge0).
\]
The complete Macaulay growth inequalities hold as well.
\end{corollary}
\begin{proof}
The $O$-sequence $d$ has embedding dimension $r-1$.
For its standard monomials, count downward edges with
coordinate multiplicity. Every selected degree-$j+1$
monomial has $j+1$ such edges, all to selected degree-$j$
monomials. Each degree-$j$ monomial has at most $r-1+j$
potential upward edge weight. Double counting gives the
inequality. The Macaulay inequalities are the classical
Hilbert-function growth conditions for an $O$-sequence.
\end{proof}

\begin{corollary}[A complete Lefschetz flag]
For every Hilbert function $h$ in the preceding theorems,
put $r=h_1$. There is a monomial realization with $r$
successive strong Lefschetz elements: the first is strong
Lefschetz on the algebra, the next on its quotient by the
first, and so on. Every iterated positive difference
$\mathcal D^k h$ is an $O$-sequence, and its embedding
dimension is $\max(r-k,0)$.
\end{corollary}
\begin{proof}
Use the previously proved recursive rising-unimodality of
ordinary arms, two-arm products, their interval windows and
positive signed single-arm kernels. Only the unimodality of
all the iterates is needed here. More generally any finite
integer sequence $h$ with $h_0=1$ whose iterates are unimodal
has such a realization. Induct on $r=h_1$. For $r=0$,
unimodality forces $h=(1)$. For $r>0$, the sequence
$g=\mathcal D h$ has $g_1=r-1$ and the same recursive
property. Realize $g$ by induction and use its standard
monomials for $\mathcal M$ in the lifting construction.
The new variable $z$ is strong Lefschetz, and the quotient
by $z$ is exactly the chosen algebra for $g$. The remaining
flag is therefore inherited from that quotient.
\end{proof}

This is an algebraic consequence of the earlier recursive
shape theorem and the classical realization criterion, not
an independent resolution of a new coefficient conjecture.
It gives Macaulay growth at every order, including
$(j+1)(\mathcal D^k h)_{j+1}\le(r-k+j)(\mathcal D^k h)_j$
for $0\le k\le r$.

This gives algebraic constraints on actual coefficients and
their positive differences, independently of a cumulative
log-concavity assertion. It does not construct a Lefschetz
algebra for a general strict one-descent core. In particular,
the SLP is not being transferred to arbitrary tensor products
or to the signed peak operation without a separate proof.
